\section{External data references and their analytical roles}\label{app:external}
The references distinguish source documentation from observations actually used in the reported calculations. ProQuest supplies the historical newspaper material and TDM Studio documents its analysis environment; the annual source exports and repository evidence map identify the received corpus. OpenAI's model documentation identifies GPT-4o mini as the model family used for the original ProQuest labels, but does not independently certify the exact snapshot or configuration of each historical call. The LOC archive supports the separate extraction-development exercise only.

FRED's CPIAUCNS, INDPRO, and UNRATE document the core macroeconomic outcomes; PAYEMS documents the alternative employment target; and USREC and the NBER chronology document recession-state definitions. The Federal Reserve Bank of St.\ Louis distributes these series through FRED, while series-specific source notes identify the original producing agencies. Shiller's monthly stock-price data are the market benchmark actually used throughout the historical evaluation. The saved workbook was downloaded on September 16, 2026; the reference-page access date below is a separate bibliographic fact.

Three additional requested FRED sources are retained as background resources: the historical Dow-Jones Industrial Stock Price Index (M1109BUSM293NNBR; \citealp{fredstock}), the ten-year Treasury yield (GS10; \citealp{fredgs10}), and the three-month Treasury bill rate (TB3MS; \citealp{fredtb3ms}). The first is an alternative historical equity series; it is not the Shiller series used in the reported market scores. The two rates could support a separately specified interest-rate comparison, but no new term-spread predictor or interest-rate forecast result is estimated here. Substituting or combining these series would require its own coverage, units, and vintage checks.

The historical EPU index is a related newspaper-based measure. The Livingston Survey, SPF, and Tealbook/Greenbook projections are candidate external forecast comparisons. None is merged into the paper's reported newspaper estimates. Their citations document the literature and the design of possible comparisons, not completed robustness tests. In particular, the public-release lag for staff forecasts must be respected when defining available information.

The August 3, 2026 retrieval dates supplied for the original reference list are retained as reported bibliographic history. Additional source checks and newly added online-resource references are dated October 21, 2026. Current resource titles and destinations are used where a supplied title or link was outdated: M1109BUSM293NNBR is labeled as a Dow-Jones index, and the Philadelphia Fed now labels its Greenbook collection Tealbook (formerly Greenbook). These bibliographic revisions do not refresh the economic data or change the saved estimates.

